% DRAFT (not yet \input in main.tex). Methodology subsection for RQ6. TODO-FINAL markers are filled from the test results.

\subsection{Guided generation with a verify-and-repair loop (RQ6)}\label{sec:method-repair}

RQ1--RQ5 measure configurations as the models write them in one attempt. RQ6 asks whether a fixed generation prompt and a generic feedback loop can turn those attempts into configurations that deploy and serve, and at what cost to security. Figure~\ref{fig:repair-arch} shows the architecture. It has three parts, and none of them uses knowledge of a particular project.

\textbf{Guided system prompt.} The neutral system prompt is replaced by a list of 17 rules (Appendix~\ref{app:guided-prompt}). The user prompt is the production prompt of RQ5, unchanged, so the guided condition differs from the production condition only in the system prompt. Every rule comes from a failure mechanism we observed on a development project (Table~\ref{tab:root-causes}), and no rule names a project, an image or a value. The rules ask the model to use only documented values, to fit hardening to what an image can run with, to give slow-starting runtimes a startup probe, to keep memory limits above what a runtime needs, to keep Service wiring consistent, and to set dependency addresses explicitly.

\textbf{Gate.} Each attempt is applied to the same cluster as in RQ1--RQ5 and passes the gate only if (i) it parses, (ii) the API server accepts every object in a server-side dry run, (iii) every pod meets Pod Security \emph{baseline}, (iv) every Service selects a pod template and targets a declared port, (v) it applies, (vi) every workload is Ready and stays Ready for 20~s within 300~s, and (vii) every Service sends traffic to a port that a selected pod actually listens on. Check (vii) reads the listening sockets of one running pod per Service from an ephemeral container in the pod's network namespace. It sends no traffic, so a NetworkPolicy cannot change its answer. We added it after the first development round, in which a configuration passed checks (i)--(vi) while its front end listened on a port that no Service used (Section~\ref{sec:res-rq6}).

\textbf{Generic feedback.} When an attempt fails the gate, the model receives one follow-up message that lists what the cluster reported: the API server's rejections, the workloads that did not become Ready, each failing container's exit code with its error log lines, warning events, and port mismatches from check (vii). It covers only objects the configuration itself created. The smoke test and the fault injection of D5/D6 are never shown to the model; they remain the independent measure of whether the configuration works. The model then returns a complete corrected configuration. The loop stops when the gate passes or after three repairs.

\textbf{Protocol.} To keep the prompt from fitting the projects it is measured on, we split the projects before any guided run. The prompt and the gate were developed on four projects (n8n, Voting App, Sock Shop, Kubernetes Dashboard) with one sample per model, and each version was kept with its results. After \TODO{TODO-FINAL: number of versions} versions the prompt was frozen and run once on the held-out projects: Grafana, Emojivoto and Online Boutique, plus three projects never used before in this study, Wordsmith, TeaStore and Headlamp. The new projects were prepared during project screening (Section~\ref{sec:method-dataset}); their human manifests pass D1--D6 on our cluster. For them we also generated neutral and production configurations (three samples per model and prompt) so that every condition has a baseline. Three models took part: Claude, Gemini and Qwen, with three samples per project. We report four conditions: neutral, production, guided (the first attempt, $r_0$) and guided with repair (the final attempt).

\textbf{Guardrails.} A loop that rewards configurations for becoming Ready could do so by removing security controls. We therefore require Pod Security baseline in the gate, and report D2 for the final attempts next to D5/D6. A repair counts as successful only by the independent smoke test, not by the gate.
